\section{Proofs for Section~\ref{sec:tasks}}\label{app:proofs-tasks}

Notation is as in Section~\ref{sec:tasks}. For a set function $f$ on the subsets of a finite set, \emph{increasing differences} means $f(\mathcal Z\cup\{i\})-f(\mathcal Z)\le f(\mathcal Z'\cup\{i\})-f(\mathcal Z')$ whenever $\mathcal Z\subseteq\mathcal Z'$ and $i\notin\mathcal Z'$. We use three elementary facts.

\begin{enumerate}
\item[(T1)] A set function has increasing differences if and only if it is supermodular. Indeed, if it has increasing differences and $B\setminus A=\{b_1,\dots,b_k\}$, then with $A_m:=A\cup\{b_1,\dots,b_m\}$ and $B_m:=(A\cap B)\cup\{b_1,\dots,b_m\}\subseteq A_m$,
\[
f(A\cup B)-f(A)=\sum_{m=1}^k[f(A_m)-f(A_{m-1})]\ge\sum_{m=1}^k[f(B_m)-f(B_{m-1})]=f(B)-f(A\cap B).
\]
The converse is the special case $A=\mathcal Z'$, $B=\mathcal Z\cup\{i\}$. Nonnegative combinations and sums of supermodular functions are supermodular.
\item[(T2)] (Weierstrass product inequality.) If $x_1,\dots,x_J$ are all nonnegative, or all lie in $[-1,0]$, then $\prod_j(1+x_j)\ge 1+\sum_jx_j$. By induction: with $P_k=\prod_{j\le k}(1+x_j)$ and $\Sigma_k=\sum_{j\le k}x_j$, $P_{k+1}=P_k(1+x_{k+1})\ge(1+\Sigma_k)(1+x_{k+1})=1+\Sigma_{k+1}+\Sigma_kx_{k+1}\ge 1+\Sigma_{k+1}$, because $1+x_{k+1}\ge 0$ and $\Sigma_k$ and $x_{k+1}$ have the same sign.
\item[(T3)] (Monotone comparative statics.) Let $f(\cdot;\mathsf t)$ be supermodular for each value of a real parameter $\mathsf t$ and have increasing differences in $(\mathcal Z,\mathsf t)$: $f(\mathcal Z';\mathsf t)-f(\mathcal Z;\mathsf t)$ is nondecreasing in $\mathsf t$ whenever $\mathcal Z\subseteq\mathcal Z'$. Let $\mathcal F\subseteq\mathcal F'$ be finite ground sets, $\mathsf t\le\mathsf t'$, $\mathcal Z$ a maximizer of $f(\cdot;\mathsf t)$ over the subsets of $\mathcal F$ and $\mathcal Z'$ a maximizer of $f(\cdot;\mathsf t')$ over the subsets of $\mathcal F'$. Then $\mathcal Z\cup\mathcal Z'$ maximizes $f(\cdot;\mathsf t')$ over the subsets of $\mathcal F'$, because it is feasible, $\mathcal Z\cap\mathcal Z'\subseteq\mathcal F$, and
\[
f(\mathcal Z\cup\mathcal Z';\mathsf t')-f(\mathcal Z';\mathsf t')\ge f(\mathcal Z;\mathsf t')-f(\mathcal Z\cap\mathcal Z';\mathsf t')\ge f(\mathcal Z;\mathsf t)-f(\mathcal Z\cap\mathcal Z';\mathsf t)\ge 0 .
\]
With $\mathsf t=\mathsf t'$ and $\mathcal F=\mathcal F'$, the union of two maximizers is a maximizer, so a largest maximizer exists; and the largest maximizer at $(\mathsf t',\mathcal F')$ contains every maximizer at $(\mathsf t,\mathcal F)$ \citep{topkis1998supermodularity,milgrom1994monotone}.
\end{enumerate}

\subsection*{Proof of Proposition~\ref{prop:tasks-modes}}

(a) By \eqref{eq:tasks-cost}, each $c_j(x;\cdot)$ is affine with slope $n_j(x)\ge 0$. The minimum of finitely many nondecreasing affine functions is nondecreasing, concave and piecewise affine. If the minimizer $x^*$ is unique at $\pi$, continuity keeps it the unique minimizer on a neighborhood, where $c^*_j=c_j(x^*;\cdot)$ has slope $n_j(x^*)$.

(b) Let $x$ be optimal at $\pi$ and $x'$ at $\pi'<\pi$. Then $c_j(x';\pi')\le c_j(x;\pi')$ and $c_j(x;\pi)\le c_j(x';\pi)$. Adding, $c_j(x';\pi')-c_j(x';\pi)\le c_j(x;\pi')-c_j(x;\pi)$, that is, $n_j(x')(\pi'-\pi)\le n_j(x)(\pi'-\pi)$, and dividing by $\pi'-\pi<0$ gives $n_j(x')\ge n_j(x)$. The intensities are $0$ for $H$, $O$ and $N$, $n^U_j$ for $U$ and $n_j>n^U_j$ for $A$. Hence once $U$ is optimal, every optimal mode at a lower price is $U$ or $A$, and once $A$ is optimal, every optimal mode at a lower price is $A$.

(c) The first statement is the definition of an optimum. If the firm does not automate, it uses the best alternative, so automating saves $\underline c_j-c_j(A)$ per unit when this is positive and nothing otherwise. The function $u\mapsto(u-c_j(A))^+$ is nondecreasing and 1-Lipschitz, and $\underline c_j\le c_j(H)$ because $H\in\mathcal M_j\setminus\{A\}$, so $0\le(c_j(H)-c_j(A))^+-(\underline c_j-c_j(A))^+\le c_j(H)-\underline c_j$. If $c_j(O)<c_j(A)<c_j(H)$, a comparison with $H$ alone recommends $A$, which costs $c_j(A)$ per unit instead of $\underline c_j\le c_j(O)$.

(d) For $x\ne A$, $n_j(x)<n_j$, so $c^0_j(A)+\pi n_j\le c^0_j(x)+\pi n_j(x)$ if and only if $\pi\le[c^0_j(x)-c^0_j(A)]/[n_j-n_j(x)]$. Automation is optimal if and only if this holds for every $x\in\mathcal M_j\setminus\{A\}$.

(e) If the alternatives have no fixed cost, $\underline c_j$ does not depend on $v_j$, and neither does $c^{A,\mathrm{var}}_j$. Automation is optimal if and only if $F^A_j/(v_jT_A)\le\underline c_j-c^{A,\mathrm{var}}_j$. With $F^A_j>0$ this requires $\underline c_j>c^{A,\mathrm{var}}_j$ and is then equivalent to $v_j\ge v^*_j$. \qed

\subsection*{Proof of Proposition~\ref{prop:tasks-curve}}

(a) With vendor prices fixed, all unit costs are affine in $\pi$ as in \eqref{eq:tasks-cost}, so by Proposition~\ref{prop:tasks-modes}(d), $A$ is optimal at $\pi$ if and only if $\pi\le\pi^*_j$ (ties are resolved in favor of $A$). Each indicator $\mathbf 1\{\pi\le\pi^*_j\}$ is a nonincreasing step function, and $S_A$ is a nonnegative combination of them; $\pi\mapsto\sum_jE_j\mathbf 1\{\pi^*_j\ge\pi\}/\sum_jE_j$ is by definition the spend-weighted survival function of crossover prices.

(b) Adding $O$ adds one alternative with $n_j(O)=0$, so the minimum in Proposition~\ref{prop:tasks-modes}(d) becomes $\pi^*_j(O)=\min\{\pi^*_j,[c_j(O)-c^0_j(A)]/n_j\}\le\pi^*_j$, and every indicator, hence $S_A$, weakly falls. Along $\pi(t)=\pi(0)\exp(-\delta_\pi t)$, a crossover price $\pi^*\in(0,\pi(0))$ is reached at $t^*=\delta_\pi^{-1}\ln[\pi(0)/\pi^*]$, so for crossovers that lie ahead the delay is $\delta_\pi^{-1}\ln[\pi^*_j/\pi^*_j(O)]$. A price path that stays positive never reaches $\pi^*_j(O)\le 0$.

(c) Since $\mathcal M_j\cap\{H,O,N\}\subseteq\mathcal M_j$, $c^*_j\le c^{\mathrm{pre}}_j$, so $\mathcal W\ge 0$. By Proposition~\ref{prop:tasks-modes}(a), each $c^*_j$ is nondecreasing and concave, so $\mathcal W$ is nonincreasing and convex, and where every optimal mode is unique its derivative is $-\sum_jv_jn_j(x^*_j(\pi))=-N(\pi)$. \qed

\subsection*{Proof of Proposition~\ref{prop:tasks-window}}

Here $c_j(O;\pi)=(1+\upsilon_O)c^V_j(\pi)+\kappa_Oc_j(H)+\zeta_{Oj}$ and $c_j(A;\pi)=c^0_j(A)+\pi n_j$.

(a) For $\pi\ge\pi^V_j$, $\pi n_j+c^{V0}_j\ge c^{VH}_j$, so $c^V_j=c^{VH}_j$ and $D_j(\pi)=(1+\upsilon_O)c^{VH}_j+\kappa_Oc_j(H)+\zeta_{Oj}-c^0_j(A)-\pi n_j$, which is strictly decreasing. For $\pi\le\pi^V_j$, $c^V_j=\pi n_j+c^{V0}_j$ and $D_j(\pi)=D_j(0)+\upsilon_O\pi n_j$, which is nondecreasing. The two expressions agree at $\pi^V_j$, so $D_j$ is continuous and attains its maximum at $\pi^V_j$. A function that is nondecreasing and then decreasing has intervals as superlevel sets; $\{D_j\ge 0\}$ is therefore an interval, nonempty if and only if $D_j(\pi^V_j)\ge 0$, in which case it contains $\pi^V_j$. Its upper end solves the decreasing branch, which gives $\pi^+_j$, and $\pi^+_j\ge\pi^V_j$ exactly when $D_j(\pi^V_j)\ge 0$.

(b) On $[0,\pi^V_j]$, $D_j(\pi)=D_j(0)+\upsilon_O\pi n_j$. If $D_j(0)\ge 0$, $D_j\ge 0$ there and $\pi^-_j=0$. If $D_j(0)<0$ and $\upsilon_O>0$, the lower end of the interval is the root $-D_j(0)/(\upsilon_On_j)$, which is at most $\pi^V_j$ exactly when $D_j(\pi^V_j)\ge 0$. If $D_j(0)<0$ and $\upsilon_O=0$, $D_j\equiv D_j(0)<0$ on $[0,\pi^V_j]$ and $D_j<D_j(\pi^V_j)<0$ beyond, so the interval is empty. For $\pi<\pi^-_j$, $D_j(\pi)<0$.

(c) $c^0_j(A)=c^{A,\mathrm{var}}_j(0)+F^A_j/(v_jT_A)$, so $D_j(0)\ge 0$ if and only if $F^A_j/(v_jT_A)\le(1+\upsilon_O)c^{V0}_j+\kappa_Oc_j(H)+\zeta_{Oj}-c^{A,\mathrm{var}}_j(0)$. If the right side, $B_j$, is positive this is $v_j\ge v^{MB}_j$; otherwise it fails for every $v_j$. Since $B_j$ does not depend on $F^A_j$, $\partial v^{MB}_j/\partial F^A_j=1/(T_AB_j)>0$, $\partial v^{MB}_j/\partial\kappa_O=-v^{MB}_jc_j(H)/B_j<0$ and $\partial v^{MB}_j/\partial\zeta_{Oj}=-v^{MB}_j/B_j<0$.

(d) Along $\pi(t)=\pi(0)\exp(-\delta_\pi t)$, $\pi(t)\in[\pi^-_j,\pi^+_j]$ for $t$ between $\delta_\pi^{-1}\ln[\pi(0)/\pi^+_j]$ and $\delta_\pi^{-1}\ln[\pi(0)/\pi^-_j]$, an interval of length $\delta_\pi^{-1}\ln(\pi^+_j/\pi^-_j)$. Part (b) gives activities that are automated in house on the window and outsourced below it, so the in-house automated share can fall as $\pi$ falls. For the AI-executed share, fix $\pi'<\pi$ and suppose activity $j$ is AI-executed at $\pi$ (ties resolved toward AI). If $A$ is optimal at $\pi$, then $c_j(A;\pi')<c_j(A;\pi)\le c_j(H)$, so $H$ is not optimal at $\pi'$; and if $\pi'\ge\pi^V_j$, then $\pi>\pi'\ge\pi^V_j$ and $D_j(\pi')>D_j(\pi)\ge 0$, so $O$ is not optimal at $\pi'$. The optimum at $\pi'$ is therefore $A$, or $O$ with $\pi'<\pi^V_j$, where the vendor uses AI. If instead $O$ is optimal at $\pi<\pi^V_j$, then at $\pi'$ the vendor still uses AI, and $c_j(O;\pi')\le c_j(O;\pi)\le c_j(H)$, so the optimum at $\pi'$ is $O$ or $A$. In both cases $j$ is AI-executed at $\pi'$. \qed

\subsection*{Proof of Proposition~\ref{prop:tasks-review}}

Write $d:=d^R_j$ and $g(\xi):=1-\xi(1+\ln(1/\xi))$, with $g(0):=1$.

(a) \emph{Selection.} Let $R(z)\in[0,1]$ be the probability that a unit with stakes $z$ is reviewed, with $\mathbb E[R(z)]=\xi$. Because errors and detection are independent of stakes, the undetected error rate is $e_j\mathbb E[1-dR(z)]=e_j(1-d\xi)$ for every such rule, so the reliability requirement and the floor depend on $\xi$ only, and the expected undetected loss is $e_jL_j\mathbb E[z(1-dR(z))]=e_jL_j(1-d\,\mathbb E[zR(z)])$. Let $z_\xi:=\ln(1/\xi)$, so that $P(z\ge z_\xi)=\xi$, and $R^*:=\mathbf 1\{z\ge z_\xi\}$. Then $\mathbb E[zR^*]-\mathbb E[zR]=\mathbb E[(z-z_\xi)(R^*-R)]+z_\xi\mathbb E[R^*-R]=\mathbb E[(z-z_\xi)(R^*-R)]\ge 0$, because the two factors have the same sign at every $z$. Reviewing the top share by stakes is therefore optimal for every $\xi$.

\emph{Cost.} With $R^*$, $\mathbb E[zR^*]=\int_{z_\xi}^\infty z\exp(-z)\,dz=(1+z_\xi)\exp(-z_\xi)=\xi(1+\ln(1/\xi))$, which gives $c^Q_j(\xi)=w_jt^R_j\xi+e_jL_j[1-d\xi(1+\ln(1/\xi))]=w_jt^R_j\xi+e_jL_j[(1-d)+dg(\xi)]$. Differentiating,
\[
c^{Q\prime}_j(\xi)=w_jt^R_j-de_jL_j\ln(1/\xi),\qquad c^{Q\prime\prime}_j(\xi)=de_jL_j/\xi>0\quad\text{on }(0,1].
\] Since $c^{Q\prime}_j\to-\infty$ as $\xi\downarrow 0$ and $c^{Q\prime}_j(1)=w_jt^R_j>0$, the unconstrained minimizer is the unique root $\xi^E_j=\exp(-w_jt^R_j/(de_jL_j))\in(0,1)$. The requirement $e_j(1-d\xi)\le\bar\varepsilon_j$ is $\xi\ge(1-\bar\varepsilon_j/e_j)/d$, so the feasible set is $[\ell_j,1]$ with $\ell_j:=\max\{(1-\bar\varepsilon_j/e_j)/d,\xi^{\min}_j\}$, nonempty when the target can be met (part (d)). By strict convexity, the minimizer over $[\ell_j,1]$ is $\xi^E_j$ if $\xi^E_j\ge\ell_j$ and $\ell_j$ otherwise, which is \eqref{eq:tasks-xi}.

(b) At $\xi^E_j$, $\ln(1/\xi^E_j)=w_jt^R_j/(de_jL_j)$, so $c^Q_j=w_jt^R_j\xi^E_j+e_jL_j[1-d\xi^E_j]-w_jt^R_j\xi^E_j=e_jL_j(1-d\xi^E_j)$. Because $c^Q_j$ is strictly convex with an interior minimizer, $c^Q_j(\xi^E_j)$ lies strictly below $\lim_{\xi\downarrow 0}c^Q_j(\xi)=e_jL_j$ and below $c^Q_j(1)=w_jt^R_j+(1-d)e_jL_j$.

(c) The economic term $\exp(-w_jt^R_j/(de_jL_j))$ rises with $e_j$, $L_j$ and $d$ and falls with $w_jt^R_j$; the reliability term $[1-(1-\bar q_j)/e_j]/d$ rises with $e_j$ and $\bar q_j$ and, where it is positive (the only case in which it can bind), falls with $d$; the floor rises with $\xi^{\min}_j$. A maximum of nondecreasing functions is nondecreasing.

\emph{Quality cost.} $g(1)=0$ and $g'(\xi)=-\ln(1/\xi)<0$ on $(0,1)$, so $g>0$ on $[0,1)$, and $c^Q_j(\xi;e)=w_jt^R_j\xi+e L_j[(1-d)+dg(\xi)]$ is strictly increasing in $e$ whenever $\xi<1$ or $d<1$. Let $e<e'$. The feasible set at $e'$ is contained in that at $e$, because $\ell_j$ is nondecreasing in $e$. Hence, if $\xi^*_j(e')<1$, $c^Q_j(e')=c^Q_j(\xi^*_j(e');e')>c^Q_j(\xi^*_j(e');e)\ge c^Q_j(e)$. Where the economic term binds strictly, $c^Q_j=e_jL_j[1-d\exp(-y/e_j)]$ locally, with $y=w_jt^R_j/(dL_j)$, and $dc^Q_j/de_j=L_j[1-d\xi^E_j(1+y/e_j)]=L_j[1-d\xi^*_j(1+\ln(1/\xi^*_j))]$. For every $\xi$, $c^Q_j(\xi;d)$ is nonincreasing in $d$, with slope $-e_jL_j\xi(1+\ln(1/\xi))\le 0$, and the feasible set expands as $d$ rises, so the minimum is nonincreasing in $d$.

\emph{Undetected rate.} $1-d\xi^*_j$ is the minimum of $1-d\xi^E_j$, $\bar\varepsilon_j/e_j$ and $1-d\xi^{\min}_j$, so
\[
\varepsilon_j=\min\bigl\{e_j(1-d\exp(-c/(de_j))),\ \bar\varepsilon_j,\ e_j(1-d\xi^{\min}_j)\bigr\},\qquad c=w_jt^R_j/L_j.
\] With $x:=c/(de_j)$, the first term has derivative $1-d(1+x)\exp(-x)\ge 1-d\ge 0$, because $(1+x)\exp(-x)\le 1$; the second is constant and the third nondecreasing, so $\varepsilon_j$ is nondecreasing in $e_j$. Since $e_j(\iota_j)$ is decreasing, $\xi^*_j$, $c^Q_j$ and $\varepsilon_j$ are nonincreasing in $\iota_j$.

(d) For $e_j\le\bar\varepsilon_j$, $e_j(1-d\xi)\le\bar\varepsilon_j$ for every $\xi\ge 0$, so the requirement never binds; for $e_j>\bar\varepsilon_j$ it excludes review shares below $(1-\bar\varepsilon_j/e_j)/d>0$. The undetected rate $e_j(1-d\xi)$ is decreasing in $\xi$, so its smallest value on $[0,1]$ is $e_j(1-d)$, and the target can be met at some share if and only if $e_j(1-d)\le\bar\varepsilon_j$, that is, $d\ge 1-\bar\varepsilon_j/e_j$; this also makes $\ell_j\le 1$, given $\xi^{\min}_j\le 1$. With $e_j(\iota)$ continuous and strictly decreasing, $\{\iota:e_j(\iota)(1-d)\le\bar\varepsilon_j\}$ is an interval $[\underline\iota_j,1]$ (possibly empty or all of $[0,1]$). As $e_j\to 0$, $(1-\bar\varepsilon_j/e_j)/d\to-\infty$ and $\xi^E_j=\exp(-c'/e_j)\to 0$ with $\exp(-c'/e_j)/e_j^k\to 0$ for every $k$ ($c'=w_jt^R_j/(dL_j)$), so $\xi^*_j\to\xi^{\min}_j$, and equals it once $\xi^E_j<\xi^{\min}_j$ if $\xi^{\min}_j>0$. Finally $0\le e_jL_j[(1-d)+dg(\xi^*_j)]\le e_jL_j\to 0$, so $c^Q_j\to\xi^{\min}_jw_jt^R_j$. \qed

\subsection*{Proof of Proposition~\ref{prop:tasks-augment}}

(a) $c_j(A;\iota)=\pi n_j+c^Q_j(e_j(\iota))+F^A_j/(v_jT_A)+(r+\delta_K)K_j(A)$. The optimal review share \eqref{eq:tasks-xi} and the cost $c^Q_j(\xi;e)$ are continuous in $e$, so $c^Q_j(e)$ is continuous; it is strictly increasing in $e$ where $\xi^*_j<1$ (Proposition~\ref{prop:tasks-review}(c)), which holds whenever $\xi^{\min}_j<1$ and $\bar q_j<1$. Composition with the continuous, strictly decreasing $e_j(\iota)$ gives the claim.

(b) Because $c_j(U)<\min\{c_j(H),c_j(O)\}$, augmentation is the best alternative to automation at every readiness level (elimination, if feasible, is assumed costlier). The function $\iota\mapsto c_j(A;\iota)-c_j(U)$ is continuous and strictly decreasing by (a), so the set where it is nonpositive is $[\iota^c_j,1]$, with $\iota^c_j$ as stated (intermediate value theorem). Automation is feasible on $[\underline\iota_j,1]$, so it is optimal exactly on $[\max\{\underline\iota_j,\iota^c_j\},1]$, and augmentation below.

(c) Automation beats augmentation only if $\pi n_j+\xi^*_jw_jt^R_j+e_jL_j[(1-d^R_j)+d^R_jg(\xi^*_j)]+F^A_j/(v_jT_A)+(r+\delta_K)K_j(A)\le w_jt^U_j+e^U_jL_j+\pi n^U_j+(r+\delta_K)K_j(U)$. Since the error term on the left is nonnegative, the stated inequality follows. With $t^R_j=t^U_j$ and $e^U_j=0$ it reads $\xi^*_j\le 1-[\pi(n_j-n^U_j)+F^A_j/(v_jT_A)+(r+\delta_K)(K_j(A)-K_j(U))]/(w_jt^U_j)$. If $\xi^{\min}_j$ exceeds the right side, then so does $\xi^*_j\ge\xi^{\min}_j$, at every readiness level.

(d) If the target can be met, the share $\xi=1$ is feasible (proof of Proposition~\ref{prop:tasks-review}(d)), so $c^Q_j(\xi^*_j)\le c^Q_j(1)=w_jt^R_j+(1-d^R_j)e_jL_j$. Hence $c_j(A)\le\pi n_j+w_jt^R_j+(1-d^R_j)e_jL_j+F^A_j/(v_jT_A)+(r+\delta_K)K_j(A)$, and with $t^R_j=t^U_j$ the right side is at most $c_j(U)=\pi n^U_j+w_jt^U_j+e^U_jL_j+(r+\delta_K)K_j(U)$ exactly under the stated inequality. \qed

\subsection*{Proof of Proposition~\ref{prop:tasks-complement}}

Write $W_p=W^u_p+W^\eta_p+W^Q_p+W^\Xi_p$, with $W^u_p(\mathcal Z)=\sum_{j\in\mathcal Z}v_j[c_j(M)-c_j(A)]$, $W^\eta_p(\mathcal Z)=\sum_ev_e[\eta_e(M,M)-\eta_e(y_i,y_j)]$, $W^Q_p(\mathcal Z)=v_pL_p[Q_p(\mathcal Z)-Q_p(\emptyset)]$ and $W^\Xi_p(\mathcal Z)=\Xi_p[\mathbf 1\{\mathcal C_p\subseteq\mathcal Z\}-\mathbf 1\{\mathcal C_p=\emptyset\}]$.

(b) Fix an edge $e=(i,j)$. If neither end is automated it contributes $0$ to $W^\eta_p(\mathcal Z)$. If only $i$ is, it contributes $v_e[\eta_e(M,M)-\eta_e(A,M)]$, which is its contribution to $W^\eta_p(\{i\})$; similarly if only $j$ is. If both are, it contributes $v_e[\eta_e(M,M)-\eta_e(A,A)]=v_e[\eta_e(M,M)-\eta_e(A,M)]+v_e[\eta_e(M,M)-\eta_e(M,A)]+v_e\eta^\times_e$, the sum of its contributions to $W^\eta_p(\{i\})$ and $W^\eta_p(\{j\})$ plus $v_e\eta^\times_e$. Summing over edges gives the decomposition.

(a) $W^u_p$ and the first sum in (b) are modular. For each edge, $\mathcal Z\mapsto\mathbf 1\{i\in\mathcal Z,j\in\mathcal Z\}$ is supermodular: if both $A$ and $B$ contain $\{i,j\}$, both sides of the supermodularity inequality equal 2; if exactly one does, the left side is at least 1; otherwise the right side is 0. With $\eta^\times_e\ge 0$, $W^\eta_p$ is supermodular by (T1). The same case check shows that $\mathcal Z\mapsto\mathbf 1\{\mathcal C_p\subseteq\mathcal Z\}$ is supermodular, so $W^\Xi_p$ is, since $\Xi_p\ge 0$. Finally $W^Q_p(\mathcal Z)=v_pL_pQ_p(\emptyset)[\prod_{j\in\mathcal Z}\hat q_j-1]$ with $\hat q_j=q_j(A)/q_j(M)$. For $\mathcal Z\subseteq\mathcal Z'$ and $i\notin\mathcal Z'$, the increments of $\Pi(\mathcal Z):=\prod_{j\in\mathcal Z}\hat q_j$ are $\Pi(\mathcal Z)(\hat q_i-1)$ and $\Pi(\mathcal Z')(\hat q_i-1)$. If all $\hat q_j\ge 1$, then $\Pi(\mathcal Z')\ge\Pi(\mathcal Z)>0$ and $\hat q_i-1\ge 0$; if all $\hat q_j\le 1$, then $0<\Pi(\mathcal Z')\le\Pi(\mathcal Z)$ and $\hat q_i-1\le 0$. Either way the increment at $\mathcal Z'$ is at least that at $\mathcal Z$, so $W^Q_p$ is supermodular by (T1), and so is the sum.

(c) Let the non-automated mode of each activity be its best alternative, with handoff costs and reliabilities that do not depend on which alternative it is. For $\mathcal Z\subseteq\mathcal Z'$, $W_p(\mathcal Z';\pi)-W_p(\mathcal Z;\pi)$ equals $\sum_{j\in\mathcal Z'\setminus\mathcal Z}v_j[\underline c_j(\pi)-c_j(A;\pi)]$ plus terms that do not depend on $\pi$, and $\underline c_j(\pi)-c_j(A;\pi)=\min_{x\ne A}\{c^0_j(x)-c^0_j(A)+\pi[n_j(x)-n_j]\}$ is nonincreasing in $\pi$ because $n_j(x)<n_j$. So $W_p$ has increasing differences in $(\mathcal Z,-\pi)$, and (T3) with parameter $-\pi$ gives the first claim; (T3) with a fixed parameter and nested ground sets of feasible activities gives the second. \qed

\subsection*{Proof of Proposition~\ref{prop:tasks-aon}}

(a) Let $n:=|\mathcal Z|$ and $\hat q$ the common ratio $q_j(A)/q_j(M)$. By the proof of Proposition~\ref{prop:tasks-complement}, $W_p(\mathcal Z)=\sum_{j\in\mathcal Z}\beta_j+\sum_{e\subseteq\mathcal Z}v_e\eta^\times_e+O(n)+\Xi_p\mathbf 1\{n=J\}$, where $\beta_j:=v_j[c_j(M)-c_j(A)]+W^\eta_p(\{j\})$ and $O(n):=v_pL_pQ_p(\emptyset)(\hat q^{\,n}-1)$. Since $\mathcal C_p=\mathcal J_p$ and $J\ge 2$, singletons do not release the overhead, so $w_j:=W_p(\{j\})=\beta_j+O(1)$, and
\[
W_p(\mathcal Z)=\sum_{j\in\mathcal Z}w_j-nO(1)+O(n)+\sum_{e\subseteq\mathcal Z}v_e\eta^\times_e+\Xi_p\mathbf 1\{n=J\}.
\]
For $n\in\{0,J\}$ the claim holds with equality. For $0<n<J$,
\[
W_p(\mathcal Z)-\frac nJW_p(\mathcal J_p)=T_1+T_2-T_3-T_4,
\]
with $T_1:=\sum_{j\in\mathcal Z}w_j-\frac nJ\sum_jw_j$, $T_2:=O(n)-\frac nJO(J)$, $T_3:=\frac nJ\sum_ev_e\eta^\times_e-\sum_{e\subseteq\mathcal Z}v_e\eta^\times_e$ and $T_4:=\frac nJ\Xi_p$ (the terms in $O(1)$ cancel). First, $T_1=\frac{J-n}J\sum_{j\in\mathcal Z}w_j-\frac nJ\sum_{j\notin\mathcal Z}w_j\le\frac{n(J-n)}J(w_{\max}-w_{\min})$. Second, $u\mapsto\hat q^{\,u}-1$ is convex on $[0,J]$ and vanishes at $0$, so $O(n)\le\frac nJO(J)$ and $T_2\le 0$. Third, by \eqref{eq:tasks-h}, $T_3\ge\frac{n(J-n)}Jh_p$. Fourth, $T_4=\frac{n(J-n)}J\cdot\frac{\Xi_p}{J-n}\ge\frac{n(J-n)}J\cdot\frac{\Xi_p}{J-1}$. Hence $W_p(\mathcal Z)-\frac nJW_p(\mathcal J_p)\le\frac{n(J-n)}J[w_{\max}-w_{\min}-h_p-\Xi_p/(J-1)]\le 0$. If $W_p(\mathcal J_p)>0$, every $\mathcal Z\ne\mathcal J_p$ has $W_p(\mathcal Z)<W_p(\mathcal J_p)$; if $W_p(\mathcal J_p)<0$, every nonempty $\mathcal Z$ has $W_p(\mathcal Z)<0=W_p(\emptyset)$.

(b) Write $E(\mathcal Z)$ for the edges with both ends in $\mathcal Z$, $r(\mathcal Z)$ for the number of connected components of the induced subgraph and $n=|\mathcal Z|$. \emph{Trees.} $|\mathcal E_p|=J-1$ and the induced subgraph is a forest, so $|E(\mathcal Z)|=n-r(\mathcal Z)\le n-1$. Then $T_3/c\ge\frac nJ(J-1)-(n-1)=\frac{J-n}J$, and the bracket in \eqref{eq:tasks-h} times $J/[n(J-n)]$ is at least $c/n\ge c/(J-1)$. Removing a leaf leaves a connected set with $n=J-1$ and $|E(\mathcal Z)|=J-2$, for which $T_3/c=(J-1)^2/J-(J-2)=1/J$ and the ratio equals $c/(J-1)$. \emph{Cycles.} $|\mathcal E_p|=J$ and a proper subset induces $r\ge 1$ paths, so $|E(\mathcal Z)|=n-r$ and $T_3/c=r\ge 1$. The ratio is at least $cJ/[n(J-n)]$, minimized by a single arc of $\lfloor J/2\rfloor$ activities, where $n(J-n)=J^2/4$ for even $J$ and $(J^2-1)/4$ for odd $J$. \emph{Complete graphs.} $|E(\mathcal Z)|=n(n-1)/2$ and $|\mathcal E_p|=J(J-1)/2$, so $T_3/c=n(J-n)/2$ and the ratio is $cJ/2$ for every $\mathcal Z$.

(c) Take a chain $1\to 2\to\cdots\to J$ with $v_e\eta^\times_e=c$, $\Xi_p=0$ and $\hat q=1$, so that $W_p(\mathcal Z)=\sum_{j\in\mathcal Z}w_j+c|E(\mathcal Z)|$, where the stand-alone values $w_j$ can be set freely through the unit savings. For $\Delta>c/(J-1)$ let $w_1=\dots=w_{J-1}=w_h:=[\Delta-(J-1)c]/J$ and $w_J=w_h-\Delta$, a dispersion of $\Delta$. Then $W_p(\mathcal J_p)=Jw_h-\Delta+(J-1)c=0$, and $W_p(\mathcal J_p\setminus\{J\})=(J-1)w_h+(J-2)c=W_p(\mathcal J_p)-w_J-c=[(J-1)\Delta-c]/J>0$. \qed

\subsection*{Proof of Proposition~\ref{prop:tasks-integration}}

(a) Evaluate \eqref{eq:tasks-process} at $\emptyset$ and $\mathcal J_p$: every unit cost, every handoff and every reliability switches from $M$ to $A$, and the overhead, paid at $\emptyset$ (as $\mathcal C_p\ne\emptyset$), is released at $\mathcal J_p$.

(b) By the decomposition in the proof of Proposition~\ref{prop:tasks-aon}(a), and since $|\mathcal C_p|\ge 2$,
\begin{gather*}
W_p(\mathcal J_p)=\sum_j\beta_j+\sum_ev_e\eta^\times_e+v_pL_pQ_p(\emptyset)\Bigl(\prod_j\hat q_j-1\Bigr)+\Xi_p,\\
\sum_jW_p(\{j\})=\sum_j\beta_j+v_pL_pQ_p(\emptyset)\sum_j(\hat q_j-1).
\end{gather*} Subtracting gives \eqref{eq:tasks-premium}. Under Proposition~\ref{prop:tasks-complement}(a), $\eta^\times_e\ge 0$, $\Xi_p\ge 0$, and the $x_j:=\hat q_j-1$ are all nonnegative or all in $[-1,0]$, so the bracket is nonnegative by (T2).

(c) By (a), $W_p(\mathcal J_p)>0$ if and only if $\sum_jv_j[c_j(A)-c_j(M)]$ is less than the sum of the other three dividends. \qed

\subsection*{Proof of Lemma~\ref{lem:tasks-oring}}

(a) Independence gives $Q_p=\prod_k(1-\varepsilon_k)$, and taking logarithms, $Q_p\ge\bar q_p$ if and only if $\sum_k-\ln(1-\varepsilon_k)\le-\ln\bar q_p$. If $\varepsilon_k\ge\varepsilon_1$ for every $k$, the sum is at least $-J\ln(1-\varepsilon_1)$, so every feasible $J$ satisfies $J\le\ln\bar q_p/\ln(1-\varepsilon_1)$, the common-error bound. With a common error, $(1-\varepsilon)^J\ge\bar q_p$ if and only if $1-\varepsilon\ge x:=\bar q_p^{1/J}\in(0,1]$. Since $1-\bar q_p=1-x^J=(1-x)\sum_{k<J}x^k\le J(1-x)$, $1-x\ge(1-\bar q_p)/J$; and $1-x=1-\exp(\ln(\bar q_p)/J)\le-\ln(\bar q_p)/J$ because $1-\exp(-y)\le y$. For $J=10$ and $\bar q_p=0.99$, $1-x=0.001005$.

(b) $Q_p=\prod_jq_j$ gives $\partial Q_p/\partial q_j=Q_p/q_j$ and $\partial^2Q_p/\partial q_i\partial q_j=Q_p/(q_iq_j)>0$ for $i\ne j$; $Q_p/q_j$ is largest where $q_j$ is smallest. \qed

\subsection*{Proof of Lemma~\ref{lem:tasks-stp}}

Let $a_j:=\ln(1/\xi_j)$, so $P(z_j<a_j)=1-\xi_j$, and a unit goes straight through if and only if $z_j<a_j$ for every $j$. Hence $\mathrm{STP}_p\le\min_jP(z_j<a_j)=\min_j(1-\xi_j)$. With common stakes, $\mathrm{STP}_p=P(z<\min_ja_j)=1-\max_j\xi_j=\min_j(1-\xi_j)$. With independent stakes, $\mathrm{STP}_p=\prod_j(1-\xi_j)$. If the vector $(z_j)$ is associated, that is, $\operatorname{Cov}(f,g)\ge 0$ for all coordinatewise nondecreasing $f,g$ of finite variance, then also $\operatorname{Cov}(f,g)=\operatorname{Cov}(-f,-g)\ge 0$ for nonincreasing $f,g$. The indicators of $\{z_1<a_1\}$ and of $\{z_j<a_j\text{ for }j\ge 2\}$ are nonincreasing, so $\mathrm{STP}_p\ge P(z_1<a_1)\,P(z_j<a_j\text{ for }j\ge 2)$. Any subvector of an associated vector is associated, so induction gives $\mathrm{STP}_p\ge\prod_j(1-\xi_j)$. \qed

\subsection*{Proof of Proposition~\ref{prop:tasks-readiness}}

(a) $W_p(\mathcal J_p;\iota)=\sum_jv_j[c_j(M)-c_j(A;\iota)]+\sum_ev_e[\eta_e(M,M)-\eta_e(A,A;\iota)]+v_pL_p[\prod_j(1-\varepsilon_j(\iota))-Q_p(\emptyset)]+\Xi_p$. The first sum is continuous and strictly increasing by Proposition~\ref{prop:tasks-augment}(a), the second continuous and nondecreasing by assumption, and the third continuous and nondecreasing because each $\varepsilon_j(\iota)=e_j(\iota)(1-\xi^*_j(\iota))$ is continuous and nonincreasing in $\iota$ (proof of Proposition~\ref{prop:tasks-review}(c)).

(b) The bound of Proposition~\ref{prop:tasks-aon}(a), $W_p(\mathcal Z;\iota)\le(|\mathcal Z|/J)W_p(\mathcal J_p;\iota)$, is a property of the value function and holds whether or not every activity is feasible. If $\iota\ge\iota^*_p$, full automation is feasible ($\iota\ge\underline\iota_p$), $W_p(\mathcal J_p;\iota)\ge 0$ by (a), and every $\mathcal Z$ has $W_p(\mathcal Z;\iota)\le W_p(\mathcal J_p;\iota)$, so full automation is optimal. If $\iota<\underline\iota_p$, it is infeasible; if $\underline\iota_p\le\iota<\iota^*_p$, then $W_p(\mathcal J_p;\iota)<0=W_p(\emptyset)$. Wherever $W_p(\mathcal J_p;\iota)<0$, every nonempty $\mathcal Z$, feasible or not, has $W_p(\mathcal Z;\iota)<0$, so nothing is automated and each activity runs in its best alternative. If $\iota^*_p>\underline\iota_p$, then $W_p(\mathcal J_p;\underline\iota_p)<0$, and by (a) $W_p(\mathcal J_p;\iota)<0$ for every $\iota<\iota^*_p$.

(c) Full automation needs each activity to be feasible. With activity-specific readiness, each term of $W_p(\mathcal J_p;\cdot)$ is nondecreasing in each $\iota_j$ (taking $\eta_e(A,A)$ nonincreasing in the readiness of both ends), and under the condition of Proposition~\ref{prop:tasks-aon}(a), full automation is optimal if and only if it is feasible and $W_p(\mathcal J_p)\ge 0$. Both conditions define upper sets, and so does their intersection. With $Q_p=\prod_jq^A(\iota_j)$, $\partial Q_p/\partial\iota_j=Q_pq^{A\prime}(\iota_j)/q^A(\iota_j)$; $q^{A\prime}$ is nonincreasing and $q^A$ is increasing and positive, so the ratio is nonincreasing in $\iota_j$ and largest at the smallest $\iota_j$.

(d) By (b), for $\iota\in[\iota^0,\iota^*_p)$ nothing is automated and the process value is $0$. The objective
\[
\max_{\mathcal Z}W_p(\mathcal Z;\iota)/r-[I_p(\iota)-I_p(\iota^0)]
\]
therefore equals $-[I_p(\iota)-I_p(\iota^0)]<0$ for $\iota\in(\iota^0,\iota^*_p)$, below its value $0$ at $\iota^0$. For $\iota\ge\iota^*_p$ the process value is $W_p(\mathcal J_p;\iota)$, which gives the investment rule. \qed

\subsection*{Proof of Proposition~\ref{prop:tasks-capital}}

(a) Capital $\Delta K_jv_j(t)$ must be in place at every date. The firm buys $\Delta K_jv_j(0)$ at date 0 and then $\Delta K_j[\dot v_j(t)+\delta_Kv_j(t)]$ a year, which is nonnegative because $v_j$ is nondecreasing (taken absolutely continuous). Let $P:=\int_0^\infty\exp(-rt)v_j(t)\,dt<\infty$. Because $v_j$ is nondecreasing, $\int_T^\infty\exp(-rt)v_j(t)\,dt\ge v_j(T)\exp(-rT)/r$, and the left side tends to 0, so $\exp(-rT)v_j(T)\to 0$. Integrating by parts, $\int_0^\infty\exp(-rt)\dot v_j(t)\,dt=-v_j(0)+rP$. The NPV is therefore
\[
\begin{aligned}
\int_0^\infty\exp(-rt)\bigl[\Delta c_jv_j-\Delta K_j(\dot v_j+\delta_Kv_j)\bigr]dt-\Delta K_jv_j(0)&=(\Delta c_j-\delta_K\Delta K_j)P-\Delta K_j[rP-v_j(0)]-\Delta K_jv_j(0)\\
&=[\Delta c_j-(r+\delta_K)\Delta K_j]P .
\end{aligned}
\]
$P=v_j/r$ for constant volume and $P=v_j(0)/(r-g)$ for $v_j(t)=v_j(0)\exp(gt)$, $g<r$. Applying the formula to any two modes shows that they are ranked by $c_j(x)+(r+\delta_K)K_j(x)$.

(b) With $v_j(t)=v_j(0)\exp(gt)$, the cash flow changes at each date by $v_j(t)[\Delta c_j-(g+\delta_K)\Delta K_j]$, and revenue grows at the same rate in a steady state. Under GAAP, depreciation $\delta_K\Delta K_jv_j(t)$ replaces capital spending. With volume proportional to revenue, the switch lowers $a$ by $v_j\Delta c_j/S$ and raises $k$ by $v_j\Delta K_j/S$; with depreciation, $m_0$ becomes $m_0-\delta_Kk$ and $\partial R^*_X/\partial k=-(g^*+\delta_K)$ (Remark~\ref{rem:assumptions}(c)), while $\partial R^*_X/\partial a=-1$ (Proposition~\ref{prop:envelope}), which gives the first-order change.

(c) The switch creates value, raises the steady-state cash margin and raises the GAAP margin if and only if $\Delta c_j$ exceeds $(r+\delta_K)\Delta K_j$, $(g+\delta_K)\Delta K_j$ and $\delta_K\Delta K_j$ respectively. If $g<r$, the first threshold is the highest, so value implies cash, with the stated band in between; if $g>r$, cash implies value, with the stated band; and cash implies GAAP because $g\ge 0$. For two modes with $c_2+(r+\delta_K)K_2=c_1+(r+\delta_K)K_1$ and $K_2>K_1$, the difference in cash cost is $[c_2+(g+\delta_K)K_2]-[c_1+(g+\delta_K)K_1]=(g-r)(K_2-K_1)$, negative if and only if $g<r$, and the difference in GAAP cost is $-r(K_2-K_1)<0$. \qed

\subsection*{Proof of Lemma~\ref{lem:tasks-success}}

(a) A product of numbers in $(0,1]$ is at most each factor, so $\varpi_p\le\varpi^0_p\le\min_j\varpi_{pj}$, and $-\ln$ turns the product into a sum. (b) For $\varpi_{pj}=1/(1+\exp(-z_j))$, $d\ln\varpi_{pj}/dz_j=1-\varpi_{pj}$, and the other factors do not depend on $z_j$; $1-\varpi_{pj}$ is largest where $\varpi_{pj}$ is smallest. As $z_i\to\infty$ for $i\ne j_0$, $\varpi_{pi}\to 1$. \qed

\subsection*{Proof of Proposition~\ref{prop:tasks-transition}}

(a) A failure is detected at an exponential time with mean $T_{\mathrm{det}}$, so it is detected during the run with probability $1-p_u$. All branches pay $F+I+C^J$ at date 0 and the parallel premium on $[0,T_{\mathrm{par}})$, worth $C^{\mathrm{par}}A_r(T_{\mathrm{par}})$. On success (probability $\varpi_p$) the saving $bW\varrho(t)$ accrues during the run and $bW$ thereafter until $T_h$; its expected present value is $\bar b_pW\mathcal S_r$, because $b$ enters linearly; and the stranded costs are paid at $T_{\mathrm{par}}$. A detected failure (probability $(1-\varpi_p)(1-p_u)$) saves nothing and, being abandoned, pays no stranded costs. An undetected failure (probability $(1-\varpi_p)p_u$) saves nothing and pays $C^{\mathrm{str}}+L_{\mathrm{roll}}$ at $T_{\mathrm{par}}$. Taking expectations gives \eqref{eq:tasks-npv}.

(b) The first statement restates (a). As $r\to 0$, $A_r(T_{\mathrm{par}})\to T_{\mathrm{par}}$, $\exp(-rT_{\mathrm{par}})\to 1$ and
\[
\mathcal S_r\to\int_0^{T_{\mathrm{par}}}\hat\varrho(t/T_{\mathrm{par}})\,dt+T_h-T_{\mathrm{par}}=T_h-(1-\bar\varrho)T_{\mathrm{par}}.
\] The undiscounted NPV is therefore $\varpi_p\bar b_pW[T_h-(1-\bar\varrho)T_{\mathrm{par}}]$ minus the bracketed costs, which is positive if and only if $T_h>T^*_h$.

(c) With $\varpi_p=b=1$ and no discounting, cumulative cash at $t\ge T_{\mathrm{par}}$ is $-(F+I+C^J)-C^{\mathrm{par}}T_{\mathrm{par}}+W\bar\varrho T_{\mathrm{par}}-C^{\mathrm{str}}+W(t-T_{\mathrm{par}})$, which is increasing in $t$ and vanishes at $t_{\mathrm{pay}}$. With $T_{\mathrm{par}}=0$ the payback date is $(F+I+C^J+C^{\mathrm{str}})/W$, so the extension is $T_{\mathrm{par}}(1-\bar\varrho+C^{\mathrm{par}}/W)$.

(d) Without discounting, and since $[\varpi_p+(1-\varpi_p)p_u]C^{\mathrm{str}}=\varpi_pC^{\mathrm{str}}+(1-\varpi_p)p_uC^{\mathrm{str}}$, the part of the NPV that depends on $T:=T_{\mathrm{par}}$ is $-f(T)$ with $f(T):=\hat cT+R\exp(-T/T_{\mathrm{det}})$ and $R:=(1-\varpi_p)(C^{\mathrm{str}}+L_{\mathrm{roll}})\ge 0$. Then $f''(T)=(R/T_{\mathrm{det}}^2)\exp(-T/T_{\mathrm{det}})\ge 0$ and $f'(T)=\hat c-(R/T_{\mathrm{det}})\exp(-T/T_{\mathrm{det}})$. If $R\le\hat cT_{\mathrm{det}}$, $f'\ge 0$ on $[0,\infty)$ and $T^*_{\mathrm{par}}=0$. Otherwise $f$ is strictly convex and $f'$ has the unique root \eqref{eq:tasks-tpar}, where $(1-\varpi_p)\exp(-T^*_{\mathrm{par}}/T_{\mathrm{det}})=(1-\varpi_p)\hat cT_{\mathrm{det}}/R=T_{\mathrm{det}}\hat c/(C^{\mathrm{str}}+L_{\mathrm{roll}})$. At an interior optimum $T^*_{\mathrm{par}}=T_{\mathrm{det}}[\ln(1-\varpi_p)+\ln(C^{\mathrm{str}}+L_{\mathrm{roll}})-\ln T_{\mathrm{det}}-\ln\hat c]$, which rises with $1-\varpi_p$, $C^{\mathrm{str}}$ and $L_{\mathrm{roll}}$ and falls with $\hat c$, hence with $\omega_{\mathrm{par}}$ and $\nu_{\mathrm{rec}}$; and $\varpi_p$ lowers $\ln(1-\varpi_p)$ and raises $\hat c$. With discounting, the first-order condition is $d\,\mathrm{NPV}_p/dT_{\mathrm{par}}=0$ applied to \eqref{eq:tasks-npv}; \texttt{code/verify\_part2\_core.py} solves it numerically for the example. \qed

\subsection*{Proof of Proposition~\ref{prop:tasks-scorepath}}

(a) On the success path the cash saving is $bW\varrho(t)$ during the run, net of the parallel premium, and $bW$ afterwards; the one-time cash charges are as stated, and with revenue and growth fixed $\Delta R_X=\Delta m$. The flow is negative exactly when $bW\varrho(t)<C^{\mathrm{par}}$, which for $\varrho(t)=t/T_{\mathrm{par}}$ is $t<T_{\mathrm{par}}C^{\mathrm{par}}/(bW)$. With $b=1$, the cumulative change in the score times $S$ is the cumulative cash flow of Proposition~\ref{prop:tasks-transition}(c).

(b) Under GAAP capital spending is replaced by depreciation, which in the steady state equals the replacement spending already netted in $W$, and the productivity dip is expensed at date 0. The recognition of the build and of the restructuring charges under the standards is Proposition~\ref{prop:acct-path} and its proof. In the simple convention, the build is amortized straight line over $T_A$ from date 0 and severance, break fees and the write-off are expensed at cutover, so after $\max\{T_{\mathrm{par}},T_A\}$ neither the ramp nor the amortization changes the margin.

(c) With $T_{\mathrm{par}}\le 1$ and a linear ramp, the flow integrates over the first year to $\int_0^{T_{\mathrm{par}}}(bWt/T_{\mathrm{par}}-C^{\mathrm{par}})\,dt+bW(1-T_{\mathrm{par}})=bW(1-T_{\mathrm{par}}/2)-C^{\mathrm{par}}T_{\mathrm{par}}$. The cash score also bears $F+I+C^J+C^{\mathrm{str}}$ in the first year; in the simple convention the GAAP score bears $C^J+C^{\mathrm{str}}+C^{\mathrm{wo}}$ and a full year of amortization $F/T_A$ (as $T_A\ge 1$). The difference is $(F+I-F/T_A-C^{\mathrm{wo}})/S$. \qed

\subsection*{Proof of Proposition~\ref{prop:tasks-pilot}}

(a) After the signal, scaling yields $\mathbb E[Y\mid\text{signal}]$ and stopping yields $0$. With a binary outcome and posterior success probability $\varpi'_p$, $\mathbb E[Y\mid\text{signal}]=\varpi'_pY_s+(1-\varpi'_p)Y_f\ge 0$ if and only if $\varpi'_p\ge|Y_f|/(Y_s+|Y_f|)$.

(b) With the pilot, the firm pays $F_{\mathrm{pil}}$ and, at $T_{\mathrm{pil}}$, takes the better of scaling and stopping, which is worth $d\,\mathbb E[(\mathbb E[Y\mid\text{signal}])^+]$ today. Without it, the firm decides now, which is worth $(\mathbb E Y)^+$. By conditional Jensen,
\[
\mathbb E\bigl[(\mathbb E[Y\mid\text{signal}])^+\bigr]-(\mathbb EY)^+\le\mathbb E[Y^+]-(\mathbb EY)^+,
\]
which is bounded as in Proposition~\ref{prop:tasks-granularity}(b). So $\mathrm{VP}\le d\min\{s/2,s^2/(4|m|)\}-F_{\mathrm{pil}}$ for the mean $m$ and standard deviation $s$ of $Y$.

(c) A revealing pilot leads to scaling exactly on success, so $\mathbb E[(\mathbb E[Y\mid\text{signal}])^+]=\varpi_pY_s$ and $\mathrm{VP}$ is as stated. On $[0,\varpi^{\mathrm{gate}}]$, where $\mathbb EY\le 0$, $\mathrm{VP}=d\varpi_pY_s-F_{\mathrm{pil}}$, which increases in $\varpi_p$. On $[\varpi^{\mathrm{gate}},1]$, $\mathrm{VP}=(1-\varpi_p)|Y_f|-(1-d)\varpi_pY_s-F_{\mathrm{pil}}$, which decreases. So $\mathrm{VP}$ is continuous and single-peaked at $\varpi^{\mathrm{gate}}$, and $\{\mathrm{VP}>0\}$ is an open interval whose ends are the roots of the two branches, $F_{\mathrm{pil}}/(dY_s)$ and $(|Y_f|-F_{\mathrm{pil}})/(|Y_f|+(1-d)Y_s)$. If $\mathrm{VP}(\varpi^{\mathrm{gate}})\le 0$, the left root is at least $\varpi^{\mathrm{gate}}$ and the right root at most $\varpi^{\mathrm{gate}}$, so the interval \eqref{eq:tasks-pilotband} is empty, as is $\{\mathrm{VP}>0\}$. With $d=1$, $\mathrm{VP}=\varpi_pY_s-F_{\mathrm{pil}}-[\varpi_pY_s-(1-\varpi_p)|Y_f|]^+=\min\{\varpi_pY_s,(1-\varpi_p)|Y_f|\}-F_{\mathrm{pil}}$. Above the band $\mathbb EY>0$ and the firm scales directly; below it $\mathbb EY<0$ and it stops. \qed

\subsection*{Proof of Proposition~\ref{prop:tasks-fullcost}}

(a) The run costs and technical debt of a one-activity process are annual amounts, so they enter per unit divided by $v_j$; with security spending chosen optimally, the incident loss and the spending together cost $C^{\mathrm{risk}}(C^{\mathrm{inc}}_j)$ a year (Proposition~\ref{prop:tasks-security}).

(b) None of these costs depends on $\pi$, so they add $c^{\mathrm{anc}}_j(x)$ to $c^0_j(x)$. By Proposition~\ref{prop:tasks-modes}(d), the crossover becomes $\min_x[c^0_j(x)+c^{\mathrm{anc}}_j(x)-c^0_j(A)-c^{\mathrm{anc}}_j(A)]/[n_j-n_j(x)]$, which falls by $\Delta c^{\mathrm{anc}}_j/[n_j-n_j(x^*)]$ if the minimizing alternative $x^*$ is unchanged. Since $c_j(A;\pi)-\underline c_j(\pi)$ is nondecreasing in $\pi$ (proof of Proposition~\ref{prop:tasks-complement}(c)), automation is optimal at some $\pi\ge 0$ if and only if it is optimal at $\pi=0$, that is, if and only if $c^0_j(A)+c^{\mathrm{anc}}_j(A)\le\underline c_j(0)$.

(c) Maintenance and revalidation do not change $c^{A,\mathrm{var}}_j$ and replace the fixed cost $F^A_j/T_A$ in Proposition~\ref{prop:tasks-modes}(e) by $F^A_j/T_A+\delta_MF^A_j+C^{\mathrm{reval}}_j/T_{\mathrm{up}}$; the ratio of the two thresholds is the stated factor. The incident cost $\Omega_j(h^f+h^uv_j)L^{\mathrm{inc}}_j$ has a part independent of $v_j$, a fixed cost, and a part proportional to $v_j$, which is a cost per unit.

(d) The term $c^{\mathrm{int}}\sum_{j\in\mathcal Z}N_j$ is modular. The term $c^{\mathrm{cpl}}\sum_{e\subseteq\mathcal Z}N_iN_j$ has the same form as $\sum_{e\subseteq\mathcal Z}v_e\eta^\times_e$ in the decomposition of Proposition~\ref{prop:tasks-complement}(b), so subtracting it replaces $v_e\eta^\times_e$ by $v_e\eta^\times_e-c^{\mathrm{cpl}}N_iN_j$, and the proofs of Propositions~\ref{prop:tasks-complement} and \ref{prop:tasks-aon} use $\eta^\times_e$ only through this sum. The risk cost is $R(x(\mathcal Z))$ with $R:=C^{\mathrm{risk}}$ concave and nondecreasing (Proposition~\ref{prop:tasks-security}(a)) and $x(\mathcal Z)=\sum_{j\in\mathcal Z}C^{\mathrm{inc}}_j$ modular and nonnegative. For $\mathcal Z\subseteq\mathcal Z'$ and $i\notin\mathcal Z'$, $x(\mathcal Z)\le x(\mathcal Z')$, so by concavity $R(x(\mathcal Z')+C^{\mathrm{inc}}_i)-R(x(\mathcal Z'))\le R(x(\mathcal Z)+C^{\mathrm{inc}}_i)-R(x(\mathcal Z))$: $R\circ x$ is submodular, and $-R\circ x$ is supermodular by (T1). If $C^{\mathrm{inc}}_j$ is the same for all $j$, $n\mapsto-R(nC^{\mathrm{inc}}_1)$ is convex and vanishes at $0$, so it satisfies the inequality used for $T_2$ in the proof of Proposition~\ref{prop:tasks-aon}(a), and that proposition holds with the risk cost included. \qed

\subsection*{Proof of Proposition~\ref{prop:tasks-buildbuy}}

(a) Annual costs are $C^{\mathrm{fix}}_b+v\pi n$ and $C^{\mathrm{fix}}_v+vc_v$, both affine in $v$; with $c_v>\pi n$ the first is smaller if and only if $v\ge v^{BB}$.

(b) $\partial(C^{\mathrm{fix}}_b+v\pi n-C^{\mathrm{fix}}_v-vc_v)/\partial\delta_M=F^A-(1+\upsilon_B)(v/V_B)F^V$, which is positive, so that a higher rate raises the cost of building relative to buying, if and only if $F^A>(1+\upsilon_B)(v/V_B)F^V$. The same computation with respect to $1/T_{\mathrm{up}}$ gives $C^{\mathrm{reval}}-(1+\upsilon_B)(v/V_B)C^{\mathrm{reval},V}$.

(c) Building is cheaper if and only if $C^{\mathrm{fix}}_b+v\pi n\le F^{\mathrm{int}}/T_A+C^{\mathrm{lock}}+vc_v$, which rearranges to $C^{\mathrm{lock}}\ge C^{\mathrm{lock}*}$. At a renewal, a price increase that costs the firm more over the coming term than switching would make it switch, so the premium per term is at most $C^{\mathrm{sw}}$, that is, at most $C^{\mathrm{sw}}/T_c$ a year without discounting. If $C^{\mathrm{sw}}/T_c<C^{\mathrm{lock}*}$, every sustainable premium is below $C^{\mathrm{lock}*}$. \qed

\subsection*{Proof of Proposition~\ref{prop:tasks-security}}

(a) Write $x:=C^{\mathrm{inc}}_p\ge 0$ and $f(C;x):=C+\mathsf g(C)x$. Because $\mathsf g$ is convex and $x\ge 0$, $f(\cdot;x)$ is convex, and $f(C;x)\ge C$ makes it coercive, so a minimizer exists. Its right derivative at $0$ is $1+\mathsf g'(0)x$, so the minimizer is positive if and only if $-\mathsf g'(0)x>1$, and an interior minimizer satisfies $1+\mathsf g'(C)x=0$. Since $\partial f/\partial C=1+\mathsf g'(C)x$ is nonincreasing in $x$ ($\mathsf g'\le 0$), $f$ has decreasing differences in $(C,x)$, and the largest minimizer is nondecreasing in $x$; exposure and the loss per incident raise $x$. $C^{\mathrm{risk}}(x)=\min_Cf(C;x)$ is a minimum of affine functions of $x$ with nonnegative slopes $\mathsf g(C)$, hence concave and nondecreasing, and $C^{\mathrm{risk}}(0)=0$. A concave function that vanishes at $0$ is subadditive on $[0,\infty)$: for $x,y>0$, $C^{\mathrm{risk}}(x)\ge\frac{x}{x+y}C^{\mathrm{risk}}(x+y)$ and $C^{\mathrm{risk}}(y)\ge\frac{y}{x+y}C^{\mathrm{risk}}(x+y)$, and adding gives $C^{\mathrm{risk}}(x)+C^{\mathrm{risk}}(y)\ge C^{\mathrm{risk}}(x+y)$. Where the minimizer is unique, Danskin's theorem gives the slope $\mathsf g(C^{\mathrm{sec}*}_p)$.

(b) With $\mathsf g(C)=\exp(-C/C_s)$, $f'(C)=1-(x/C_s)\exp(-C/C_s)$. If $x\le C_s$, $f'\ge 0$ on $[0,\infty)$, so $C^{\mathrm{sec}*}_p=0$ and $C^{\mathrm{risk}}=x$. If $x>C_s$, the root is $C^{\mathrm{sec}*}_p=C_s\ln(x/C_s)$, where $\mathsf g(C^{\mathrm{sec}*}_p)x=C_s$ and $C^{\mathrm{risk}}=C_s\ln(x/C_s)+C_s$. With $u:=x/C_s>1$, $C^{\mathrm{sec}*}_p/x=\ln(u)/u$, which is maximized at $u=e$ with value $1/e$.

(c) The exclusions by data class and tail loss are constraints set by risk appetite. Otherwise, automating $j$ in addition to $\mathcal Z$ raises the risk cost by $C^{\mathrm{risk}}(x(\mathcal Z)+C^{\mathrm{inc}}_j)-C^{\mathrm{risk}}(x(\mathcal Z))\le C^{\mathrm{risk}}(C^{\mathrm{inc}}_j)$ by subadditivity, and for a stand-alone activity the test is $W_j>C^{\mathrm{risk}}(C^{\mathrm{inc}}_j)$, which under (b) is as stated. Because $\Omega_j(U)\le\Omega_j(A)$, $C^{\mathrm{inc}}_j(U)\le C^{\mathrm{inc}}_j(A)$ and, $C^{\mathrm{risk}}$ being nondecreasing, the risk cost of $U$ is no larger. \qed

\subsection*{Proof of Proposition~\ref{prop:tasks-secreview}}

(a) In the first-come, first-served queue with Poisson arrivals at rate $\nu$ and exponential service at rate $\Phi:=\Phi_{\mathrm{sec}}>\nu$, the stationary number in the system is geometric, $P(N=n)=(1-\varrho_q)\varrho_q^n$ with $\varrho_q:=\nu/\Phi$, and by the property that Poisson arrivals see time averages an arriving deployment finds this distribution. By memorylessness its time in the system is the sum of $N+1$ independent exponential times with rate $\Phi$, so its Laplace transform is
\[
\mathbb E\exp(-sT_{\mathrm{sec}})=\sum_{n\ge 0}(1-\varrho_q)\varrho_q^n\Bigl(\frac{\Phi}{\Phi+s}\Bigr)^{n+1}=\frac{(1-\varrho_q)\Phi}{\Phi+s-\varrho_q\Phi}=\frac{\Phi-\nu}{\Phi-\nu+s},
\]
the transform of an exponential distribution with rate $\Phi-\nu$. Hence $\mathbb ET_{\mathrm{sec}}=1/(\Phi-\nu)$, and a deployment that waits $T_{\mathrm{sec}}$ keeps $V_{\mathrm{dep}}\exp(-r_dT_{\mathrm{sec}})$ in expectation, so it loses $V_{\mathrm{dep}}[1-(\Phi-\nu)/(\Phi-\nu+r_d)]=V_{\mathrm{dep}}r_d/(\Phi-\nu+r_d)$.

(b) Multiplying by the arrival rate gives $D_{\mathrm{sec}}$, and differentiating gives \eqref{eq:tasks-secvalue}.

(c) $D_{\mathrm{sec}}(\Phi)+p_{\mathrm{sec}}\Phi$ is convex in $\Phi$ on $(\nu,\infty)$, and its derivative vanishes where $(\Phi-\nu+r_d)^2=\nu V_{\mathrm{dep}}r_d/p_{\mathrm{sec}}$.

(d) If deployments arrive at rate $\nu\ge\Phi_{\mathrm{sec}}$, the number waiting grows without bound, so the long-run deployment rate is at most $\Phi_{\mathrm{sec}}$. With a decision latency $\ell_{\mathrm{dec}}$ followed by an independent review, the total delay is $\ell_{\mathrm{dec}}+T_{\mathrm{sec}}$, and $\mathbb E\exp(-r_d(\ell_{\mathrm{dec}}+T_{\mathrm{sec}}))=\exp(-r_d\ell_{\mathrm{dec}})(\Phi_{\mathrm{sec}}-\nu)/(\Phi_{\mathrm{sec}}-\nu+r_d)$ by (a). \qed

\subsection*{Proof of Proposition~\ref{prop:tasks-ancscale}}

(a) The classification follows Assumption~\ref{ass:tasks-ancillary} and Proposition~\ref{prop:tasks-fullcost}(c); dividing a fixed amount by more processes or more volume lowers it per unit.

(b) A cell treated as one process bears one per-process fixed cost $\bar\phi$ instead of $\sum_p\phi_p$. With $P$ equal processes, true ancillary cost is $P\phi+U$, with $U$ the per-unit part, and the coarse estimate is $\phi+U$; the gap $(P-1)\phi$ is the share $(1-1/P)P\phi/(P\phi+U)=(1-1/P)s_{pp}$ of the truth. By Jensen's inequality, $\sum_pC^{\mathrm{td}}(N_p)\ge PC^{\mathrm{td}}(\bar N)$ for convex $C^{\mathrm{td}}$; for $C^{\mathrm{td}}(N)=bN+cN^2/2$ (plus any affine term) the gap is $\frac c2[\sum_pN_p^2-P\bar N^2]=\frac c2P\operatorname{Var}(N)$, with $\operatorname{Var}$ the variance across processes. Subadditivity of $C^{\mathrm{risk}}$ (Proposition~\ref{prop:tasks-security}(a)), applied repeatedly, gives $\sum_pC^{\mathrm{risk}}(C^{\mathrm{inc}}_p)\ge C^{\mathrm{risk}}(\sum_pC^{\mathrm{inc}}_p)$. \qed

\subsection*{Proof of Proposition~\ref{prop:tasks-halluc}}

(a) As in the proof of Proposition~\ref{prop:tasks-review}, errors are independent of stakes, so reviewing the top share by stakes minimizes the expected undetected loss for each $\xi$. An unreviewed unit with stakes $z$ loses $L_jz$ with probability $\bar e_j$; a reviewed one loses it with probability $\bar e_j-\tilde e_j$. With $\mathbb E[z\mathbf 1\{z\ge z_\xi\}]=\xi(1+\ln(1/\xi))$ and $\mathbb Ez=1$, the expected loss is $L_j[\bar e_j-\tilde e_j\xi(1+\ln(1/\xi))]=\tilde e_jL_j[1-\xi(1+\ln(1/\xi))]+(\bar e_j-\tilde e_j)L_j$. The last term does not depend on $\xi$, so the minimization in Proposition~\ref{prop:tasks-review} applies with $e_jL_j$ replaced by $\tilde e_jL_j$. Equivalently, the loss is $\bar e_jL_j[1-(\tilde e_j/\bar e_j)\xi(1+\ln(1/\xi))]$, the quality cost of Proposition~\ref{prop:tasks-review} with raw error rate $\bar e_j$ and detection rate $\tilde e_j/\bar e_j$.

(b) The undetected rate is $\bar e_j(1-\xi)+(\bar e_j-\tilde e_j)\xi=\bar e_j-\tilde e_j\xi$, which is at most $\bar\varepsilon_j$ if and only if $\xi\ge(\bar e_j-\bar\varepsilon_j)/\tilde e_j$; some $\xi\le 1$ satisfies this if and only if $\bar e_j-\tilde e_j\le\bar\varepsilon_j$, which is the condition $d^R_j\ge 1-\bar\varepsilon_j/\bar e_j$ of Proposition~\ref{prop:tasks-review}(d) with $d^R_j=\tilde e_j/\bar e_j$. The floor $(\bar e_j-\tilde e_j)L_j=[e^o_j(1-d^o_j)+e^h_j(1-d^h_j)]L_j$ is linear in $d^h_j$ with slope $-e^h_jL_j$. \qed

\subsection*{Proof of Proposition~\ref{prop:tasks-bias}}

(a) Reviewing the top share $\xi$ of $u\sim\mathrm{Exp}(1)$ means reviewing $u\ge a:=\ln(1/\xi)$, and $\mathbb E[u\mid u\ge a]=1+a$ by memorylessness, so the prevalence among reviewed units is $e_j(1+a)$, decreasing in $\xi$. Errors in the reviewed pool number $e_j\xi(1+a)$ per unit, and each is caught with probability $d^R(\bar e(\xi))$.

(b) With $d^R(\bar e)=d^R_0(\bar e/\bar e_0)^{\epsilon_d}$, $\mathsf k(\xi)=\mathrm{const}\cdot\xi(1+a)^{1+\epsilon_d}$, and since $da/d\xi=-1/\xi$,
\[
\frac{d\ln\mathsf k}{d\xi}=\frac1\xi-\frac{1+\epsilon_d}{\xi(1+a)}=\frac1\xi\Bigl[1-\frac{1+\epsilon_d}{1+\ln(1/\xi)}\Bigr],
\]
which is positive if and only if $\ln(1/\xi)>\epsilon_d$, that is, $\xi<\exp(-\epsilon_d)$. Any $\xi>\exp(-\epsilon_d)$ costs more review than $\exp(-\epsilon_d)$ and catches fewer errors, so it is dominated.

(c) At fixed $\xi$, $\mathsf k\propto e_j^{1+\epsilon_d}$, so $\partial\mathsf k/\partial e_j=(1+\epsilon_d)\mathsf k/e_j$, and $\partial(e_j-\mathsf k)/\partial e_j=1-(1+\epsilon_d)\mathsf k/e_j$, which is negative if and only if $\mathsf k/e_j>1/(1+\epsilon_d)$.

(d) The undetected rate is $e_j-\mathsf k(\xi)$, and by (b) $\mathsf k$ attains its maximum on $(0,1]$ at $\xi_b:=\exp(-\epsilon_d)$, so the target can be met if and only if $e_j-\mathsf k(\xi_b)\le\bar\varepsilon_j$. At $\xi_b$, $\ln(1/\xi_b)=\epsilon_d$, so $\xi_b(1+\ln(1/\xi_b))=\exp(-\epsilon_d)(1+\epsilon_d)$, the prevalence is $e_j(1+\epsilon_d)$ and $\mathsf k(\xi_b)=e_j\exp(-\epsilon_d)(1+\epsilon_d)\,d^R(e_j(1+\epsilon_d))$. Rearranging $\mathsf k(\xi_b)\ge e_j-\bar\varepsilon_j$ gives the condition. Finally $\exp(\epsilon_d)>1+\epsilon_d$ for $\epsilon_d>0$. \qed

\subsection*{Proof of Proposition~\ref{prop:tasks-autonomy}}

(a) For $W=\min\{\mathsf Z,\bar{\mathsf Z}\}$, $\mathbb EW=\int_0^{\bar{\mathsf Z}}P(\mathsf Z>x)\,dx=1+\int_1^{\bar{\mathsf Z}}x^{-\mathsf p}dx$, which is $1+(1-\bar{\mathsf Z}^{1-\mathsf p})/(\mathsf p-1)$ for $\mathsf p\ne 1$ and $1+\ln\bar{\mathsf Z}$ for $\mathsf p=1$. As $\bar{\mathsf Z}\to\infty$ it tends to $1+1/(\mathsf p-1)=\mathsf p/(\mathsf p-1)$ if $\mathsf p>1$ and diverges if $\mathsf p\le 1$. A level with $e^m_j(1-d^m_j)L^m_j>0$ and an infinite expected loss cannot minimize $c_j$ when some level has finite cost (drafting has none). For the tail, the quantile of $\mathsf Z$ at level $u$ is $(1-u)^{-1/\mathsf p}$, so for $\mathsf p>1$ the conditional value at risk at level $\mathsf q$ is $(1-\mathsf q)^{-1}\int_{\mathsf q}^1(1-u)^{-1/\mathsf p}du=\mathsf p(1-\mathsf q)^{-1/\mathsf p}/(\mathsf p-1)$; scaling by $L^m_j$ gives the stated value, and a capped loss never exceeds $L^{\mathrm{cap}}_j$, so neither does its conditional value at risk.

(b) The levels form a finite chain, and $\mathsf a+1$ is better than $\mathsf a$ if and only if $c_j(\mathsf a+1)<c_j(\mathsf a)$, which is the stated comparison.

(c) $c_j(\mathsf a+1)-c_j(\mathsf a)=w_j[t_j(\mathsf a+1)-t_j(\mathsf a)]+\Delta c^Q_j(\mathsf a)+\theta_m[\hat c^m_j(\mathsf a+1)-\hat c^m_j(\mathsf a)]$ is nondecreasing in $\theta_m$ because $\hat c^m_j$ is nondecreasing, so $c_j$ has increasing differences in $(\mathsf a,\theta_m)$ and, by (T3) applied to $-c_j$ on the chain with parameter $-\theta_m$, the largest minimizer is nonincreasing in $\theta_m$. The term $w_j[t_j(\mathsf a+1)-t_j(\mathsf a)]$ is nonincreasing in $w_j$ because $t_j$ is nonincreasing, which gives the claim for the wage. With a common cap, $\mathbb E\min\{\mathsf Z,\bar{\mathsf Z}\}$ is a common factor of all levels and is decreasing in $\mathsf p$, so it acts as such a $\theta_m$.

(d) The risk appetite is a constraint on the feasible levels. If level 3 violates it, the feasible levels are at most 2, and approval is required; if level 2 also violates it, so do levels 0 and 1 only if they act, which they do not, so automation, which requires action, is infeasible and the activity is done in another mode. \qed

\subsection*{Proof of Proposition~\ref{prop:tasks-pipeline}}

(a) Independence gives $(1-\varepsilon)^J$ per run and $(1-\varepsilon)^{Jk}$ for $k$ runs. The condition $(1-\varepsilon)^{Jk}\ge\bar q_p$ is $J\le\ln\bar q_p/[k\ln(1-\varepsilon)]$, and for small $1-\bar q_p$ and $\varepsilon$ both logarithms are close to $-(1-\bar q_p)$ and $-\varepsilon$.

(b) A segment of $J_s$ steps fails with probability $1-(1-\varepsilon)^{J_s}$, and a failure escapes the checkpoint with probability $1-d^c$. With $J/J_s$ independent segments, the condition $\{1-[1-(1-\varepsilon)^{J_s}](1-d^c)\}^{J/J_s}\ge\bar q_p$ gives the bound. For small $\varepsilon$, $1-(1-\varepsilon)^{J_s}\approx J_s\varepsilon$, so the bound is about $\ln\bar q_p/[-\varepsilon(1-d^c)]$, which is $J_{\max}/(1-d^c)$ with $k=1$.

(c) An error first occurs at step $i$ with probability $\varepsilon(1-\varepsilon)^{i-1}$ and then costs $L^c(J-i+1)$. With $x=1-\varepsilon$, $\sum_{i=1}^J(J-i+1)x^{i-1}=\sum_{m=1}^J(1-x^m)/(1-x)=[J-x(1-x^J)/(1-x)]/(1-x)$, and multiplying by $L^c\varepsilon$ gives the closed form. Expanding, $(1-\varepsilon)(1-(1-\varepsilon)^J)/\varepsilon=J-\varepsilon J(J+1)/2+O(\varepsilon^2)$. The budget condition $L^c\varepsilon J(J+1)/2\le\bar L$ gives $J\approx\sqrt{2\bar L/(\varepsilon L^c)}$ for large $J$.

(d) Under a constant hazard, $\prod_k2^{-d_k/\mathcal H}=2^{-\sum_kd_k/\mathcal H}$. Under the logistic form, a step of length $x$ succeeds with probability $p(x)=[1+(x/\mathcal H)^{1/\eta_r}]^{-1}$, so $-\ln p(x)=\phi(x):=\ln(1+(x/\mathcal H)^{1/\eta_r})$. For $\eta_r\ge 1$, $x\mapsto(x/\mathcal H)^{1/\eta_r}$ is concave and nondecreasing and $u\mapsto\ln(1+u)$ is strictly concave and increasing, so $\phi$ is strictly concave on $[0,\infty)$ with $\phi(0)=0$, hence strictly subadditive: $\phi(x)+\phi(y)>\phi(x+y)$ for $x,y>0$. Splitting a step of length $x+y$ into independent steps of lengths $x$ and $y$ therefore raises $-\ln$ of the pipeline's success probability. \qed

\subsection*{Proof of Proposition~\ref{prop:tasks-cerr}}

(a) By Proposition~\ref{prop:tasks-halluc}(a), the expected undetected loss at the review share $\xi^*_j$ is the sum of the first two terms, and Proposition~\ref{prop:tasks-autonomy} gives the misbehavior term at the optimal level. The first term is nonnegative because $1-\xi(1+\ln(1/\xi))\ge 0$ on $(0,1]$ (proof of Proposition~\ref{prop:tasks-review}(c)), and the other two do not depend on $\xi$. (b) Summing over automated activities and dividing by revenue gives a share of revenue that, like $c_c$ and $f$, is a cost that remains when $c_h=0$ and $a=a_{\min}$, so it is subtracted inside the after-tax margin term of Corollary~\ref{cor:ceiling}. (c) With no review, every error and hallucination is undetected and $\xi^*_j=0$, so the loss is $\bar e_jL_j$ plus the misbehavior term. With a model reviewer, an error is caught only if the reviewer does not share it (probability $1-\rho^{\mathrm{cm}}_j$) and then catches it (probability $d^A_j$), independently of the stakes, so Proposition~\ref{prop:tasks-review} applies with $d^R_j=(1-\rho^{\mathrm{cm}}_j)d^A_j$; the loss in (a) is nonincreasing in the catch rate, so human rates above this one understate it. \qed

\subsection*{Proof of Corollary~\ref{cor:tasks-failure-success}}

Each step's undetected rate is at least its review-proof rate (Proposition~\ref{prop:tasks-halluc}(b)). If one such rate exceeds $1-\bar q_p$, the product of the step reliabilities is below $\bar q_p$ on every outcome, and $\varpi_p=0$. Adding a step multiplies the product by a factor in $(0,1]$, and raising a step's hallucination or misbehavior rate, or its autonomy level (which raises $e^m_j$), raises its undetected rate on every outcome, so in both cases the event that the target is met shrinks and its probability does not rise. A checkpoint lowers undetected failure (Proposition~\ref{prop:tasks-pipeline}(b)) and so enlarges the event. \qed

\subsection*{Proof of Proposition~\ref{prop:tasks-granularity}}

(a) Among decisions measurable with respect to $\mathcal I_T$, automating on an event $D$ yields $\mathbb E[Y_j\mathbf 1_D]=\mathbb E[\mathbb E[Y_j\mid\mathcal I_T]\mathbf 1_D]\le\mathbb E[(\mathbb E[Y_j\mid\mathcal I_T])^+]$, with equality for $D=\{\mathbb E[Y_j\mid\mathcal I_T]>0\}$. For $\mathcal I_T\subseteq\mathcal I_{T'}$, $\mathbb E[Y_j\mid\mathcal I_T]=\mathbb E[\mathbb E[Y_j\mid\mathcal I_{T'}]\mid\mathcal I_T]$, and conditional Jensen for the convex function $u\mapsto u^+$ gives $(\mathbb E[Y_j\mid\mathcal I_T])^+\le\mathbb E[(\mathbb E[Y_j\mid\mathcal I_{T'}])^+\mid\mathcal I_T]$. Taking expectations and summing, $\mathcal V_T\le\mathcal V_{T'}$; full information gives $\sum_j\mathbb E[Y_j^+]$.

(b) Work conditionally on $\mathcal I_T$, with $Y=Y_j$, $m=m_j$ and $s=s_j$. If $m\ge 0$, then $\ell=\mathbb E[Y^+]-m=\mathbb E[Y^-]$, and $Y^-\le(m-Y)^+$, so $\ell\le\mathbb E[(m-Y)^+]=\tfrac12\mathbb E|Y-m|\le s/2$. For $m>0$ also $Y^-\le(Y-m)^2/(4m)$: at $y\ge 0$ the left side is $0$, and at $y<0$, $(y-m)^2+4my=(y+m)^2\ge 0$. So $\ell\le s^2/(4m)$. If $m<0$, $\ell=\mathbb E[Y^+]$, $Y^+\le(Y-m)^+$ gives $\ell\le s/2$, and $Y^+\le(Y-m)^2/(4|m|)$ (at $y>0$, $(y+|m|)^2\ge 4y|m|$) gives $\ell\le s^2/(4|m|)$.

(c) For $Y\sim N(m,s^2)$, $\mathbb E[Y^+]=m\mathsf N(m/s)+s\mathsf n(m/s)$. If $m\ge 0$, $\ell=s\mathsf n(u)-m[1-\mathsf N(u)]=s\mathsf L(u)$ with $u=m/s$; if $m<0$, $\ell=-|m|[1-\mathsf N(u)]+s\mathsf n(u)=s\mathsf L(u)$ with $u=|m|/s$, by the symmetry of $\mathsf n$. Differentiating, $\mathsf L'(u)=-u\mathsf n(u)-[1-\mathsf N(u)]+u\mathsf n(u)=-[1-\mathsf N(u)]<0$, and $\mathsf L(0)=\mathsf n(0)=1/\sqrt{2\pi}\approx 0.399$; since $\mathsf L$ decreases to $0$, it is positive. For the bound let $\mathsf M(u):=(1+u^2)[1-\mathsf N(u)]-u\mathsf n(u)$. Then $\mathsf M'(u)=2u[1-\mathsf N(u)]-(1+u^2)\mathsf n(u)-\mathsf n(u)+u^2\mathsf n(u)=-2\mathsf L(u)<0$ and $\mathsf M(u)\to 0$ as $u\to\infty$, so $\mathsf M\ge 0$, that is, $u[1-\mathsf N(u)]\ge u^2\mathsf n(u)/(1+u^2)$, which gives $\mathsf L(u)\le\mathsf n(u)/(1+u^2)$ for $u\ge 0$. For the value of a finer tier, if $Y$ and the new data are jointly normal given $\mathcal I_T$, then $\tilde m:=\mathbb E[Y\mid\mathcal I_{T'}]$ is normal with mean $m$ and, by the law of total variance with the constant Gaussian posterior variance $s'^2$, variance $s^2-s'^2=\hat s^2$. The value of the tier is $\mathbb E[\tilde m^+]-m^+=\hat s\mathsf L(|m|/\hat s)$ by the formula just proved.

(d) By independence, $\mathbb E[Y_j\mid\mathcal I_T,E_j]=E_jm_\delta-F^A_j/T_A$. If $F^A_j=0$, its sign is that of $m_\delta$ whenever $E_j>0$, so the decision does not change, and the value of observing $E_j$ is $\mathbb E[(E_jm_\delta)^+]-(m_\delta\mathbb EE_j)^+=0$. If $F^A_j>0$ and $m_\delta>0$, the value is $\mathbb E[(E_jm_\delta-F^A_j/T_A)^+]-(m_\delta\mathbb EE_j-F^A_j/T_A)^+=m_\delta\{\mathbb E[(E_j-E^*_j)^+]-(\mathbb EE_j-E^*_j)^+\}$. If $m_\delta\le 0$, the activity is never automated and the value is $0$.

(e) Activity-level information yields $\sum_jY_j^+$; one mode for the cell yields $(\sum_jY_j)^+=(\sum_jY_j^+-\sum_jY_j^-)^+$. The difference is $\sum_jY_j^-$ if $\sum_jY_j^+\ge\sum_jY_j^-$ and $\sum_jY_j^+$ otherwise, that is, the smaller of the two. \qed

\subsection*{Proof of Proposition~\ref{prop:tasks-aggregate}}

(a) The changes in $c_h$, $c_c$, $a$ and $k$ are accounting identities per dollar of revenue, with builds treated as fixed costs as in Remark~\ref{rem:assumptions}(a). If growth spend buys a fixed bundle of growth activities whose cost falls by the proportion $\delta_\sigma$, a spend share $\sigma$ buys $\sigma/(1-\delta_\sigma)$ units of the original bundle, so \eqref{eq:growth} becomes $g=G+\lambda(1-\delta_\sigma)^{-\alpha}\sigma^{\alpha}h(s)$. Proposition~\ref{prop:ai} and Lemma~\ref{lem:prices} give $\gamma'=-(1-\varphi)u'/p$, and here $\Delta u/p=\Delta c_h+\Delta c_c$. When every rival makes the savings, the price falls from 1 to $p_1=1-\bar\varphi^{\rm ind}S_u$ per unit of today's output; fees stay the share $f$ of price and the other costs, $C_0-S_u$ per unit, are fixed per unit, so the margin intercept is $1-f-(C_0-S_u)/p_1$. Its derivative in $S_u$ at $S_u=0$ is $1-\bar\varphi^{\rm ind}C_0$, and since $d\ln p/d\ln u=(dp/du)(u/p)=\bar\varphi^{\rm ind}C_0$ at that point, this is Lemma~\ref{lem:prices} with the elasticity $\varphi=\bar\varphi^{\rm ind}C_0$. The floor $c^{\mathrm{floor}}_j$ is the cost of the cheapest feasible mode once inference is free and raw errors vanish, which is the limit of full automation. The expected-misstatement part of the floor is Proposition~\ref{prop:acct-amin}(c); it is not part of $c_{\mathrm{err}}$, whose losses $L_j$ exclude misstatements.

(b) Equation \eqref{eq:Rstar} holds at fixed $s$ with spend re-optimized, and with depreciation $m_0$ is replaced by $m_0-\delta_Kk$ (Remark~\ref{rem:assumptions}(c)). The first-order change follows from Proposition~\ref{prop:envelope}: $\partial R^*_X/\partial m_0=1$, $\partial R^*_X/\partial G=X-k$, $\partial R^*_X/\partial k=-(g^*+\delta_K)$ and $\partial R^*_X/\partial\lambda=(X-k)\Delta^*/\lambda$, with $\Delta\ln\lambda=-\alpha\ln(1-\delta_\sigma)=\alpha\delta_\sigma+O(\delta_\sigma^2)$ and $\alpha(X-k)\Delta^*=\sigma^*$ by \eqref{eq:sigstar}. The exact effect of $\delta_\sigma$ is Corollary~\ref{cor:priorities}(ii) with $1+x_\lambda=(1-\delta_\sigma)^{-\alpha}$. The function $f_\sigma(\delta):=(1-\delta)^{-\alpha/(1-\alpha)}-1$ is convex with $f_\sigma(0)=0$ and $f_\sigma'(0)=\alpha/(1-\alpha)$, so $f_\sigma(\delta)\ge\alpha\delta/(1-\alpha)$, and $\Pi=(1-\alpha)(X-k)\Delta^*=\frac{1-\alpha}{\alpha}\sigma^*$ gives $\Pi f_\sigma(\delta_\sigma)\ge\sigma^*\delta_\sigma$.

(c) Adding $Gv_0$ to both sides of \eqref{eq:gordon} gives $rv_0=m_0+(v_0-k)G+C_\alpha[\lambda(v_0-k)_+]^{1/(1-\alpha)}$, which is $R^*_X$ at $X=v_0$ and $h=1$ by \eqref{eq:Rstar} (or Proposition~\ref{prop:score}(c) if $v_0\le k$). Let $\mathcal G(v,\mathsf t):=R^*_v(\mathsf t)-rv$, where $\mathsf t$ is any one of the parameters. By Proposition~\ref{prop:envelope}, $\mathcal G$ is continuously differentiable with $\partial\mathcal G/\partial v=g^\circ-r<0$ at $v_0$ (Proposition~\ref{prop:dynamic}(d)), so by the implicit function theorem $dv_0/d\mathsf t=\partial_{\mathsf t}R^*_X|_{X=v_0}/(r-g^\circ)$. Since $v(s)/s\to v_0$ as $s\to 0$ (Proposition~\ref{prop:dynamic}(d)), firm value is approximately $Sv_0$ at small share, and one-time costs reduce it one for one. \qed

\subsection*{Proof of Corollary~\ref{cor:tasks-weights}}

Each statement is a coefficient of the first-order expansion in Proposition~\ref{prop:tasks-aggregate}(b), and the effect on the multiple follows from part (c). \qed

\subsection*{Proof of Proposition~\ref{prop:tasks-breakpoints}}

Items (1) and (2) are Proposition~\ref{prop:tasks-modes}(d) and (e). For (3), $c_j(A)\le\underline c_j$ is $c^Q_j(e_j)\le\underline c_j-\pi n_j-F^A_j/(v_jT_A)-(r+\delta_K)K_j(A)$. By Proposition~\ref{prop:tasks-review}, $c^Q_j(e)$ is continuous and strictly increasing on $(0,1]$ where $\xi^*_j<1$, with limit $\xi^{\min}_jw_jt^R_j$ as $e\to 0$, so when the right side lies between that limit and $c^Q_j(1)$ the intermediate value theorem gives a unique $e^*_j$, and the inequality holds if and only if $e_j\le e^*_j$; feasibility requires $d^R_j\ge 1-\bar\varepsilon_j/e_j$ by Proposition~\ref{prop:tasks-review}(d). Item (4) restates Propositions~\ref{prop:tasks-augment} and \ref{prop:tasks-readiness}. For (5), the feasible activities $\mathcal F(\mathcal H)=\{j:\bar d_j\le\mathcal H\}$ form a nested family that changes only when $\mathcal H$ crosses some $\bar d_j$, and $W_p$ does not depend on $\mathcal H$; by (T3), the largest maximizer over the subsets of $\mathcal F(\mathcal H)$ is nondecreasing in $\mathcal H$ and constant between consecutive $\bar d_j$. If $\mathcal H<\bar d_p$, some activity is infeasible, every feasible set is a proper subset of $\mathcal J_p$ with nonpositive value, and not automating is optimal. If $\mathcal H\ge\bar d_p$, every activity is feasible and $W_p(\mathcal J_p)>0\ge W_p(\mathcal Z)$ for every $\mathcal Z\ne\mathcal J_p$, so full automation is the unique optimum. Item (6) is Proposition~\ref{prop:tasks-transition}(b), item (7) is Propositions~\ref{prop:tasks-fullcost}(b) and \ref{prop:tasks-security}(c), item (8) is Proposition~\ref{prop:tasks-secreview}(d), and item (9) restates Propositions~\ref{prop:tasks-halluc}(b), \ref{prop:tasks-autonomy}(d) and \ref{prop:tasks-pipeline}. The automated share is a nonnegative combination of indicators, each a step function of the variable in question, and under (5) process $p$ contributes $E_p/\sum_jE_j$ from $\bar d_p$ on. \qed
